\chapter{Cardiac Catheterization and Coronary Angiography}

\section*{What Is This Test?}
Cardiac catheterization involves threading a thin tube through a blood vessel to your heart. Contrast dye is injected to visualize coronary arteries on X-ray. It's the definitive test for coronary artery blockages and allows immediate treatment (stenting).

\section*{Alternative Names}
Cardiac cath, coronary angiography, coronary angiogram, heart catheterization, left heart catheterization


\section*{Principle}
Cardiac catheterization uses a catheter inserted via the femoral or radial artery, advanced to the coronary ostia under fluoroscopic guidance. Iodinated contrast is injected, and cineangiography records the contrast flow through coronary arteries. Intravascular ultrasound and fractional flow reserve (FFR) can be performed simultaneously.

\section*{History}
Cardiac catheterization was pioneered by Werner Forssmann in 1929, who inserted a catheter into his own heart. The technique was refined by André Cournand and Dickinson Richards in the 1940s, earning all three the Nobel Prize.

\section*{Why This Test Is Ordered}
Your doctor orders cardiac catheterization when you have a positive stress test, when you are having a heart attack, when your heart failure needs evaluation, or when your doctor needs to measure pressures inside your heart. It answers the question: which of your coronary arteries are blocked, how severe are the blockages, and how well is your heart pumping?

\section*{Is This a Primary or Secondary Test}
A secondary test. Performed after non-invasive tests suggest the need for direct visualization of the coronary arteries.

\section*{How This Test Is Performed}
A thin tube (catheter) is inserted into an artery in your wrist or groin and guided to your heart using X-ray imaging. Dye is injected into your coronary arteries, and X-ray video shows exactly where blockages are. If a blockage is found, a balloon may be inflated to open the artery (angioplasty) and a stent (small metal mesh tube) may be placed to keep it open. You are awake but sedated during the procedure. It takes one to two hours.

\section*{What Preparation Is Needed}
Preparation varies with urgency, sedation, and institutional protocol. The team will tell you whether fasting is needed, review your medicines, and arrange any necessary tests. Do not stop a prescribed medicine unless the treating team specifically instructs you to do so.

\section*{Contraindications and Precautions}
Cardiac catheterization is an invasive procedure with specific contraindications including uncorrected bleeding disorders, severe kidney dysfunction, active systemic infection, and severe allergy to contrast dye (relative). Preparation varies with urgency and institutional protocol.

\section*{Risks}
Cardiac catheterization is an invasive procedure. Risks include: bleeding or bruising at the catheter insertion site (wrist or groin), allergic reaction to contrast dye, kidney injury from contrast dye, damage to the blood vessel, heart rhythm abnormalities, heart attack, stroke, and (very rarely) death. The overall complication rate is low (less than 1--2%).

\section*{What Happens After the Test}
You will lie flat for several hours to prevent bleeding at the catheter insertion site. If the procedure was performed through your wrist, recovery is faster. You may be kept overnight in the hospital. Results are available immediately during the procedure.

\section*{Understanding Your Results}

\section*{Normal Results}
A normal result shows no significant coronary artery stenosis (<50% narrowing), normal coronary flow, and normal heart pressures.

\section*{Follow-up Tests}
Echocardiogram, cardiac MRI, or coronary artery bypass grafting (CABG) evaluation.

\section*{References}
\begin{enumerate}
  \item MedlinePlus. Cardiac Catheterization and Coronary Angiography. U.S. National Library of Medicine; 2025.
  \item Current Medical Diagnosis and Treatment. McGraw-Hill; 2025.
\end{enumerate}

\clearpage
